Order-driven markets are trading venues organised around \emph{limit orders}.
A limit order is the fundamental action that market participants in order-driven markets can perform.
It is represented by a 4-tuple $(t,q,p,d)$,
where $t$ denotes time, $q$ denotes size, $p$ denotes price, and $d$ denotes direction.
We say that a market participant \emph{posts} a limit order $(t,q,p,d)$
if at time $t$ she submits to the exchange her commitment to buy ($d=1$) or sell ($d=-1$)
the amount $q$ at the price $p$.
The price $p$ is interpreted as the highest price at which she is committed to buy if $d=1$,
or the lowest price at which she is committed to sell if $d=-1$.
By regulation, $p$ must be an integer multiple of some fixed $\tickSizeOfLOB>0$.
Market participants have also the possibility to withdraw their commitments to trade,
by \emph{cancelling} a previously submitted limit order.

A trading epoch in an order-driven market is the collection of all the limit orders
submitted by market participants within some time interval.
Limit orders cannot be submitted simultaneously,
so that a limit order is identified by its timestamp $t$.
We express this mathematically by saying that a trading epoch is the graph of a function
from a subset of the positive half-line to the space
$\lbrace (q,p,d): \, q\geq0, p\geq0, d=-1,1\rbrace$.

When a market participant submits a limit order $(t,q,p,d)$,
a matching algorithm run by the trading platform looks for possible counterparts to her order:
if other participants' commitments to trade exist to (partially) fulfill her order
(at a price not worse than $p$ from her point of view),
then her order is (partially) cleared.
The fraction of her order that is cleared disappears from the market,
and it is said to be \emph{executed};
the fraction that could not be fulfilled is recorded in the \emph{limit order book},
waiting for counterparts to trade with.

A limit order $(t,q,p,d)$ in a trading epoch is said to be \emph{active} (or outstanding) at time $u$
if $t\leq u$ and by time $u$ the order has neither been executed nor been cancelled.

The search for possible trading counterparts for the matching of the incoming order $(t,q,p,d)$
happens among the active orders with timestamp $s<t$ and opposite direction $-d$.
We say that the active order $(s,\rho,\pi,-d)$ matches the incoming order $(t,q,p,d)$
if
$\pi d \leq p d$;
in this case, the two orders are executed one against the other,
and $\rho \wedge q$ units of shares are transacted at the price $\pi$ per unit of share.
After the match, the matched orders are replaced with the orders
$(s,\tilde{\rho}, \pi, -d)$ and $(t,\tilde{q}, p, d)$,
where $\tilde{\rho} = \rho - \rho \wedge q$
and $\tilde{q} = q - \rho \wedge q$.
At least one of them has null size.
If $\tilde{\rho} = 0$, i.e. $q \geq \rho$,
the order with time stamp $s$ ceases to be active and it is deleted from its queue.
If $\tilde{q} = 0$, i.e. $q\leq \rho$,
then the order with time stamp $t$ (i.e. the incoming order) is fully executed,
and the search for possible matching counterparts stops.
If instead $\tilde{q} > 0$, i.e. $q>\rho$,
then the search continues among active orders with opposite direction $-d$.

The search follows the priority given by the ordering
\begin{equation}
\label{eq.orderingOfLimitOrders}
 (s,\rho,\pi,-d)<(s\derivative,\rho\derivative,\pi\derivative,-d)
 \qquad
 \text{ if and only if }
 \qquad
 (d\pi<d\pi\derivative)\vee((\pi=\pi\derivative)\wedge(s<s\derivative)),
\end{equation}
whereby $(s,\rho,\pi,-d)$ is given priority over $(s\derivative,\rho\derivative,\pi\derivative,-d)$
if $ d\pi<d\pi\derivative$ or, in case $\pi=\pi\derivative$, $s<s\derivative$.
This means that if both $(s,\rho,\pi,-d)$ and $(s\derivative,\rho\derivative,\pi\derivative,-d)$
are potential matches to the incoming order $(t, q,p,d)$
(i.e.
$\pi d \leq p d$ and
$\pi\derivative d \leq p d$),
then
$(t, q, p, d)$ is executed first against $(s, \rho, \pi, -d)$
and only its outstanding part $(t, \tilde{q}, p, d)$ after this transaction
is executed against $(s\derivative, \rho\derivative, \pi\derivative, -d)$.

\begin{lemma}
\label{lemma.orderingOfOrders}
Let $E$ be a trading epoch in an order-driven market.
Then, equation \eqref{eq.orderingOfLimitOrders} defines a total order in
$E\cap \lbrace (s,\rho,\pi,-1): s\geq0, \rho\geq0,\pi\geq0\rbrace$,
and it defines a total order in
$E\cap \lbrace (s,\rho,\pi,1): s\geq0, \rho\geq0,\pi\geq0\rbrace$.
\end{lemma}
\begin{proof}
 Antisymmetry follows from the fact that $E$ is a graph of a function defined on a subset of $\R$.
 Transitivity and connexity are clear.
\end{proof}
Lemma \ref{lemma.orderingOfOrders} justifies the following definition.
\begin{defi}
\label{def.orderQueue}
 Let $E$ be a trading epoch in an order-driven market.
 The ask order queue $\askOrderQueue_t$ at time $t$ is defined as
 \begin{equation*}
  \askOrderQueue_t := \lbrace (s,q,p,-1)\in E : \, (s,q,p,-1) \text{ is active at time } t \rbrace.
 \end{equation*}
 Similarly, the bid order queue $\bidOrderQueue_t$ at time $t$ is defined as
 \begin{equation*}
  \bidOrderQueue_t := \lbrace (s,q,p,+1) \in E : \, (s,q,p,+1) \text{ is active at time } t \rbrace.
 \end{equation*}
\end{defi}

A limit order book is a grid of equally spaced prices at which active limit orders sit.
The space between consecutive grid nodes is called the tick size of the LOB,
which we denote by $\tickSizeOfLOB$.
As mentioned above,
all orders submitted to the exchange must have prices that are integer multiples of the tick size.
Prices are increasing from left to right.
At every node of the grid,
outstanding limit orders to buy or sell at the corresponding price are collected;
these limit orders are also said to be queuing there.
Spontaneously -- i.e.
by market forces -- such orders organise in a way that
buy offers will be displaced on the left (the so-called \emph{bid side})
and sell offers will be on the right (the so-called \emph{ask side}).

Indeed,
if there were a buy (respectively, sell) limit order to the right (left) of some sell (buy) limit orders,
or at the same node,
then the matching algorithm would have matched them, clearing them out from the order book.

Therefore, the whole configuration of a limit order book at time $t$ is given
when the following variables are specified:
\begin{enumerate}[label={\emph{\roman{*}}.}, ref={\emph{\roman{*}}.}]
\item\label{item.bestAskPrice} the best ask price $\bestAskPrice_{t}$, i.e.
the lowest price at which one can find sell offers active at time $t$;
\item\label{item.bestBidPrice} the best bid price $\bestBidPrice_t$, i.e.
the highest price at which one can find buy offers active at time $t$;
\item\label{item.askVolume} the volume $\nthBestAskVolume[i]_t$ of ask offers at price
$\nthBestAskPrice[i]_t= \bestAskPrice_t + (i-1)\tickSizeOfLOB$, for $i=1,2,\dots$, i.e.
the quantity
\begin{equation*}
 \nthBestAskVolume[i]_t = \sum_{\lbrace(s,q,\nthBestAskPrice[i]_t,-1) \in \askOrderQueue_t : \quad s\leq t\rbrace} q,
\end{equation*}
where the sum is over all the active sell limit orders submitted by time $t$
for a price of $\bestAskPrice_t + (i-1)\tickSizeOfLOB$;
\item\label{item.bidVolume} the volume $\nthBestBidVolume[i]_t$ of bid offers at price
$\nthBestBidPrice[i]_t=\bestBidPrice_t - (i-1)\tickSizeOfLOB$, for $i=1,2,\dots$, i.e.
the quantity
\begin{equation*}
 \nthBestBidVolume[i]_t = \sum_{\lbrace (s,q,\nthBestBidPrice[i]_t,1) \in \bidOrderQueue_t: \quad s\leq t\rbrace} q,
\end{equation*}
where the sum is over all the active buy limit orders submitted by time $t$
for a price of $\bestBidPrice_t - (i-1)\tickSizeOfLOB$.
\end{enumerate}
Points \ref{item.bestAskPrice}-\ref{item.bidVolume} above describe how the configuration at time $t$
of the limit order book is reconstructed from the ask queue $\askOrderQueue_t$
and the bid queue $\bidOrderQueue_t$.

\begin{remark}
\label{remark.volumeIsQueueSize}
In points \ref{item.askVolume} and \ref{item.bidVolume} the term 'volume' is used
to denote the sum of the sizes of orders queuing at a specified price level.
In particular,
this means that it is measured in number of shares.
Often in the literature this sum is also referred to as 'size' of the queue.
\end{remark}

The whole configuration of a limit order book at time $t$ is described by the variables
$(\bestAskPrice_t, \bestBidPrice_t, \lbrace (\nthBestAskVolume[i]_t, \nthBestBidVolume[i]_t) : \, \, i=1,2,\dots\rbrace)$.
Apart from these variables,
other derived quantities are useful to assess the properties of an order book.
These are the spread, the mid-price and the volume imbalance (or queue imbalance).

The spread at time $t$ is the distance $\abs{\bestAskPrice_t - \bestBidPrice_t}$
between best ask price and best bid price,
which we denote by $\LOBspread_t$.
The mid-price at time $t$ is the mid point in between $\bestBidPrice_t$ and $\bestAskPrice_{t}$,
which we denote by $\midPrice_t$ and is given by
$\midPrice_t = (\bestAskPrice_t + \bestBidPrice_{t})/{2}$.
Finally, the $n$-levels volume imbalance (or queue imbalance) at time $t$,
denoted $\volumeImbalance^{n}_t$,
is the normalised excess of limit orders on the first $n$ levels of the bid side
compared to the limit orders on the first $n$ levels of the ask side, namely
\begin{equation}
\label{eq.volumeImbalance}
\volumeImbalance^{n}_t =
\frac{\sum_{i\leq n}\nthBestBidVolume[i]_t - \sum_{i\leq n}\nthBestAskVolume[i]_t}{\sum_{i\leq n}\nthBestBidVolume[i]_t + \sum_{i\leq n}\nthBestAskVolume[i]_t},
\end{equation}
where $\sum_{i\leq n}\nthBestBidVolume[i]_t$
(respectively, $\sum_{i\leq n}\nthBestAskVolume[i]_t$)
is the cumulative volumes on the first $n$ bid (ask) levels.
The queue imbalance $\volumeImbalance^n$ is widely accepted as a reliable signal
for the next mid-price move (\cite{CDJ18enh}):
when it is close to -1 the mid-price will likely decrease,
and when it is close to +1 it will likely increase.

When a limit order comes to the market, it modifies the previous configuration of the order book.
Proposition \ref{prop.lobUpdate} gives the updating rule for the variables
$(\bestAskPrice_t, \bestBidPrice_t, \lbrace (\nthBestAskVolume[i]_t, \nthBestBidVolume[i]_t) : \, \, i=1,2,\dots\rbrace)$.
\begin{prop}
\label{prop.lobUpdate}
 Let $(t,q,p,-1)$ be a sell limit order that arrives to the market at time $t$.
 Let $(\bestAskPrice_{t-}, \bestBidPrice_{t-},\lbrace (\nthBestAskVolume[i]_{t-}, \nthBestBidVolume[i]_{t-}) : \, \, i=1,2,\dots\rbrace )$
 represent the limit order book immediately before $t$.
 Then, prices and volumes are updated according to the following formulae:
 \begin{equation}
 \label{eq.lobUpdateSellOrder}
  \begin{split}
   \bestBidPrice_t =& \nthBestBidPrice[1+N]_{t-};
   \\
   \nthBestBidVolume[k]_{t}
   =& \nthBestBidVolume[k+N]_{t-}
   - \left( q^{k+N-1}-q^{k+N}\right), \qquad k=1,2,\dots ;
   \\
   \bestAskPrice_t = & \bestAskPrice_{t-}
   - \max\left(0,\bestAskPrice_{t-} - p\right) \indicator{q^{\infty} > 0} ;
   \\
   \nthBestAskVolume[k]_t = &  \nthBestAskVolume[k+\delta\bestAskPrice_t/\tickSizeOfLOB]_{t-}
   +q^{\infty}\indicator{p=\bestAskPrice_t + (k-1)\tickSizeOfLOB}, \qquad k=1,2,\dots ;
  \end{split}
 \end{equation}
 where $\delta\bestAskPrice_t = \bestAskPrice_t - \bestAskPrice_{t-}$
 denotes the (non-positive) jump of the best ask price;
 $N:=\max(0,N_v \wedge N_p -1)$,
 where $N_v := \inf\lbrace n\geq 0 : \, q<\sum_{i=1}^{n} \nthBestBidVolume[i]_{t-}\rbrace$
 and $N_p:=\inf\lbrace n\geq 1: \, \nthBestBidPrice[n]_{t-}< p \rbrace$; and
 \begin{equation*}
  \begin{split}
   q^{i}:=& \max\left( 0,
    q-\sum_{k=1}^{i\wedge N_v \wedge (N_p - 1)} \nthBestBidVolume[k]_{t-}
   \right)
   \\
   =& \max\left( 0,
   q-\sum_{1 \leq k\leq i} \nthBestBidVolume[k]_{t-}\indicator{\nthBestBidPrice[k]_{t-} \geq p}
   \right)
  \end{split}
 \end{equation*}
 is the outstanding quantity to be executed after the first $i$ levels have been consumed.
 For a buy limit order $(t,q,p,+1)$, the updating rules are similar, with opposite side and direction.
\end{prop}
\begin{proof}
 The processing of the sell limit order $(t,q,p,-1)$ is equivalent to the processing of the stream
 $(t,\nthBestBidVolume[1]_{t-},\nthBestBidPrice[1]_{t-},-1), \dots, (t,\nthBestBidVolume[N]_{t-},\nthBestBidPrice[N]_{t-},-1), (t,q^{N}, p, -1)$,
 where the priority is from left to right.
 Hence, we are reduced to discuss the following three cases.
 \begin{enumerate}[label={\emph{\roman{*}}.}]
  \item No match with any buy order.
  This happens if $p>\bestBidPrice_{t-}$.
  In this case therefore,
  $\bestBidPrice_t = \bestBidPrice_{t-}$ and $\nthBestBidVolume[k]_{t} = \nthBestBidVolume[k]_{t-}$
  are the correct update rules on the bid side, because no modification happens.
  With $N_p=1$ (so that $N=0$) and $q^i=q$ for all $i\geq 0$
  these updates are retrieved from the formulae in the statement.
  As for the ask side,
  $\bestAskPrice_t$ must be updated to the minimum between $\bestAskPrice_{t-}$ and $p$,
  and the stated updating rule yields this.
  Finally, we have that $-\delta\bestAskPrice_t/\tickSizeOfLOB$ is the number of levels
  by which the ask price is improved by the incoming order at time $t$.
  Hence, if indeed there is an improvement (i.e. $p<\bestAskPrice_{t-}$),
  we have that $\lbrace \nthBestAskVolume[k]_{t}: \, k=1,2,\dots\rbrace$
  represents the volumes of $\lbrace \nthBestAskVolume[k]_{t-}: \, k=1,2,\dots\rbrace$
  where the index $k$ is shifted by $-\delta\bestAskPrice_t/\tickSizeOfLOB$,
  and the best ask volume at time $t$ is the size $q$ of the incoming order.
  If instead there is no improvement (i.e. $p\geq \bestAskPrice_{t-}$),
  the only volume $\nthBestAskVolume[k]_{t-}$ that is modified
  is the one corresponding to the price $p=\bestAskPrice_t + (k-1)\tickSizeOfLOB$,
  and such modification is the addition of $q$.
  \item Complete execution at the first bid level.
  This happens if $p\leq\bestBidPrice_{t-}$ and $q<\nthBestBidVolume[1]_{t-}$.
  In this case the ask side is not altered,
  and the stated updating rules yield that because $q^{\infty}=0$.
  As for the bid side, we must have
  $\bestBidPrice_t=\bestBidPrice_{t-}$,
  $\nthBestBidVolume[1]_t = \nthBestBidVolume[1]_{t-}- q$,
  and $\nthBestBidVolume[k]_t = \nthBestBidVolume[k]_{t-}$ for $k\geq 2$.
  Since in this case $N_v=1$ and $N_p\geq 2 $, $q^0 = q$ and $q^i = 0$ for all $i\geq 1$
  the stated updating rules encompass this case as well.
  \item Partial execution at first bid level.
  This happens if $p=\bestBidPrice_{t-}$ and $q\geq \nthBestBidVolume[1]_{t-}$.
  In this case, we must check that the stated updating rule yield
  $\bestBidPrice_t = \nthBestBidPrice[2]_{t-}$,
  $\nthBestBidVolume[k]_t = \nthBestBidVolume[k+1]_{t-}$ for all $k\geq 1$, on the bid side.
  As $N_v\geq 2$, $N_p=2$, $q^0 = q$, and $q^i = q-\nthBestBidVolume[1]_{t-}$ for all $i\geq 1$,
  this is indeed checked.
  Finally, on the ask side if $q=\nthBestBidVolume[1]_{t-}$ there are not modifications at all.
  If instead $q>\nthBestBidVolume[1]_{t-}$ the updating rules must yield
  $\bestAskPrice_t = p$, $\nthBestAskVolume[1]_t = q-\nthBestBidVolume[1]_{t-}$,
  and the volume index $k$ shifted by the number of levels in the spread at time $t-$.
  Since with $q>\nthBestBidVolume[1]_{t-}$ it holds $q^{\infty} = q-\nthBestBidVolume[1]_{t-}$
  and $\delta\bestAskPrice_t = -\LOBspread_{t-}$ this is confirmed.
 \end{enumerate}
\end{proof}

The arrival of a limit order to the market triggers two events:
one is the consumption of the liquidity capable to (partially) service the limit order,
the other is the addition of the non-executed part of the limit order
to the appropriate queue in the order book.
Proposition \ref{prop.decompositionOfLimitOrder} states the terms of this decomposition.

\begin{prop}
\label{prop.decompositionOfLimitOrder}
 Given the configuration
 $(\bestAskPrice_{t-}, \bestBidPrice_{t-},\lbrace (\nthBestAskVolume[i]_{t-}, \nthBestBidVolume[i]_{t-}) : \, \, i=1,2,\dots\rbrace )$
 of the limit order book immediately before time $t$,
 processing the sell limit order $(t,q,p,-1)$ is equivalent to processing the pair of orders
 $[(t,\marketOrderSize,0,-1),(t,q-\marketOrderSize,p,-1) ]$,
 with the former having priority over the latter, and where
 \begin{equation*}
  \marketOrderSize:= \min\left(q, \sum_{i\geq 1} \nthBestBidVolume[i]_{t-} \indicator{\nthBestBidPrice[i]_{t-}\geq p}\right).
 \end{equation*}
Similarly, processing the buy limit order $(t,q,p,+1)$ is equivalent to processing the pair of orders
$[(t,\marketOrderSize,\infty,+1),(t,q-\marketOrderSize,p,+1) ]$, where
 \begin{equation*}
  \marketOrderSize:= \min\left (q, \sum_{i\geq 1} \nthBestAskVolume[i]_{t-} \indicator{\nthBestAskPrice[i]_{t-}\leq  p}\right).
 \end{equation*}
\end{prop}
\begin{proof}
 Let $(t,q,p,-1)$ be a sell limit order.
 Let $N_v$, $N_p$ and $q^i$ be as in Proposition \ref{prop.lobUpdate}.
 Let $N_v^M$, $N_p^M$ and $\marketOrderSize^{i}$ be the corresponding quantities
 for the order $(t,\marketOrderSize,0,-1)$.
 Notice that $N_v^M = N_v^M \wedge N_p^M$.
 Processing $(t,\marketOrderSize,0,-1)$ does not affect the ask side
 because $\marketOrderSize^{\infty} = 0$;
 moreover the effects on the bid side of processing $(t,\marketOrderSize,0,-1)$
 are the same as those of $(t,q,p,-1)$,
 because $\inf \lbrace n\geq 0 : \, \marketOrderSize < \sum_{i=1}^{n} \nthBestBidVolume[i]_{t-} \rbrace = N_v \wedge N_p$
 and $q^{k+N-1} - q^{k+N} = \marketOrderSize^{k+N_v^M - 1} - \marketOrderSize^{k+N_v^M}$.
 Therefore, after $(t,\marketOrderSize,0,-1)$ has been processed,
 the bid side is the same as the bid side after the processing of $(t,q,p,-1)$.

 Processing $(t,q-\marketOrderSize,p,-1)$ after $(t,\marketOrderSize,0,-1)$
 does not alter the bid side because either $\bestBidPrice_{t-} < p$ or $q-\marketOrderSize=0$.
 Moreover, $(t,q-\marketOrderSize,p,-1)$ has the same effects on the ask side
 as those of $(t,q,p,-1)$ because $q-\marketOrderSize = q^{\infty}$.

 The proof of the decomposition of a buy limit order is mutatis mutandis the same.
\end{proof}

The first component $(t,\marketOrderSize,0,-1)$
in the decomposition $[(t,\marketOrderSize,0,-1),(t,q-\marketOrderSize,p,-1) ]$
of the sell limit order $(t,q,p,-1)$ is called \emph{sell market order}.
Notice that in the 4-tuple $(t,\marketOrderSize,0,-1)$, the price $p$ is set to zero,
and this guarantees immediate execution:
the sale of the amount $\marketOrderSize$ is instantaneously matched
with outstanding buy limit orders on the bid side,
and no fraction of $\marketOrderSize$ is put into the queue.
On the contrary,
the second component $(t,q-\marketOrderSize,p,-1)$ specifies exactly that part of $(t,q,p,-1)$
which will be queued.
Similarly, $(t,\marketOrderSize,\infty,+1)$ represents the market order component of
the buy limit order $(t,q,p,+1)$,
and it is referred to as \emph{buy market order}.
The price $p=\infty$ is a way to express the fact that a counterpart
for the purchase of the amount $\marketOrderSize$
will instantaneously be found on the ask side of the order book.
In the following, the term 'market order' (either buy or sell) will be referring
to the first component in the decomposition of a limit order
with non-zero market order size $\marketOrderSize >0$.\footnote{
In some trading venues, traders can actually submit market orders,
i.e. buy or sell orders that are executed without price constraints,
at least as long as offers with the opposite direction exist.
Even in such cases, we keep our convention of referring to the fraction of a limit order
that is executed upon submission as market order.}

A sell market order $(t,\marketOrderSize,0,-1)$ is said to 'walk the book'
if $\marketOrderSize > \nthBestBidVolume[1]_{t-}$;
this implies $N\geq 1 $ in the notation of Proposition \ref{prop.lobUpdate},
the converse failing when the order exactly clears the best level.
Similarly, a buy market order walks the book
if its size is larger than the volume of offers on the first ask level.

Given a time window $\timeWindow$,
the evolution in time of the limit order book
$\lbrace (\bestAskPrice_t, \bestBidPrice_t,\lbrace (\nthBestAskVolume[i]_t, \nthBestBidVolume[i]_t) : \, \, i=1,2,\dots\rbrace): \quad 0\leq t \leq \timeHorizon \rbrace$
results from the history of all limit order submissions,
cancellations and executions that happened in $\timeWindow$.
If every limit order in $\lbrace (t,q,p,d): \, 0\leq t\leq \timeHorizon\rbrace$
is decomposed as per in Proposition \ref{prop.decompositionOfLimitOrder},
then the seller-initiated trades that happened in $\timeWindow$ are
$\lbrace (t,\marketOrderSize,0,-1):\, \, \marketOrderSize>0, \, 0\leq t\leq \timeHorizon \rbrace$,
and the buyer-initiated trades that happened in $\timeWindow$ are
$\lbrace (t,\marketOrderSize,\infty,+1):\, \, \marketOrderSize>0, \, 0\leq t\leq \timeHorizon \rbrace$.
Hence, in the following we will identify trades with market orders
of non-zero size $\marketOrderSize>0$.\footnote{
Our identification is kept in place in trading venues
where traders can actually submit market orders alongside limit orders.
In these trading venues the actual trigger of a trade might be a limit order $(t,q,p,d)$
with non-trivial price specification $0<p<\infty$ or a genuine market order.
Because of the asserted equivalence in Proposition \ref{prop.decompositionOfLimitOrder},
referring to both these triggers as ``market orders''
does not alter the evolution of the limit order book.
}

\subsection{PnL equations of trade execution}
\label{sec.pnlOfTradeExecution}

Assume that
we want to execute an order of size
$\volume$
at time $t$.
We can send a market order to the order book,
which guarantees immediate execution at the cost of paying the bid-ask spread.
More precisely,
we can execute the (signed) order $\volume$
at time $t$
via a market order,
and this changes our cash account by
\begin{equation}
\label{eq.cashMarketOrder}
	\Delta \cashAccount_t^{\text{market order}}
	=
	-\midPrice_{t} \volume
	\underbrace{ -\frac{\bestAskPrice_t - \bestBidPrice_t}{2}\abs{\volume}}_{\text{we pay the spread}}
	,
\end{equation}
where
$\midPrice_t$,
$\bestAskPrice_t$,
$\bestBidPrice_t$
are respectively the mid price, the ask price and the bid price at time $t$,
and $\abs{\volume}$ is the unsigned size of our order.
The quantity $\bestAskPrice_t - \bestBidPrice_t$ is the spread at time $t$,
which we denote $\LOBspread_t$.

Instead,
if we send a limit order at time $t$,
the execution of our order will not be immediate,
but it will happen at some time $u>t$.
We re-label the time when the limit order is sent as $t-1$,
and the time when the order is hit as $t$.
In this way,
the formula for the change in our cash account at execution time $t$ reads
\begin{equation}
\label{eq.cashLimitOrder}
	\Delta \cashAccount_t^{\text{limit order}}
	=
	-\midPrice_{t-1}\volume
+
	\underbrace{\frac{\bestAskPrice_{t-1} - \bestBidPrice_{t-1}}{2}\abs{\volume}}_{\text{we earn the spread}}
\end{equation}
Notice that,
whereas in the equation for the market order only the time $t$ appeared,
in the case of limit orders
the cash account is updated
when the limit order is hit at time $t$,
and it changes as a function of the bid and ask price as the (stale)
time $t-1$ (when the order was sent).
This is the source of the distinction between
the change in portfolio value after execution of a market order
and
the change in portfolio value after execution of a limit order.
The former reads
\begin{equation}
\label{eq.wealthMarketOrder}
	\begin{split}
		\Delta \wealth_t^{\text{market order}}
		&=
		\underbrace{
			\inventory_{t-1}\Delta \midPrice_t}_{\text{PnL from position held at time} t-1}
		\underbrace{
			-\midPrice_{t}\volume
		- \frac{\LOBspread_t}{2}\abs{\volume}
	}_{\Delta \cashAccount_t^{\text{market order}}}
		\underbrace{
			+\midPrice_t \volume
		}_{\text{mark-to-market of acquired position}}
		\\
		&=
		\inventory_{t-1} \Delta \midPrice_t
		- \frac{\LOBspread_t}{2}\abs{\volume}
	\end{split}
\end{equation}
The latter reads
\begin{equation}
\label{eq.wealthLimitOrder}
	\begin{split}
		\Delta \wealth_t^{\text{limit order}}
		&=
		\underbrace{
			\inventory_{t-1}\Delta \midPrice_t}_{\text{PnL from position held at time} t-1}
		\underbrace{
			-\midPrice_{t-1}\volume
		+ \frac{\LOBspread_{t-1}}{2}\abs{\volume}
	}_{\Delta \cashAccount_t^{\text{limit order}}}
		\underbrace{
			+\midPrice_t \volume
		}_{\text{mark-to-market of acquired position}}
		\\
		&=
		\inventory_{t-1} \Delta \midPrice_t
		+\volume \Delta \midPrice_t
		+ \frac{\LOBspread_{t-1}}{2}\abs{\volume}
	\end{split}
\end{equation}

Hence,
when deciding whether to execute via market orders or via limit orders,
we need to assess
the size of the spread relative to the risk
of an adverse price move
between the time at which we send our order
and
the time at which our order is hit.
Notice that
we do not have control on the latter;
it will instead depend on other market participants,
who might or might not have
a) lower latency between order generation, submission and acknowledgement,
b) superior information on intraday price swings.
In other words,
limit orders incur the risk of
\emph{adverse selection}
by faster or more informed market participants.
The decision between market orders and limit orders
rests
on whether
or not
the bid-ask spread adequately compensates for this risk.
